\documentclass[10pt,a4paper]{article}
\usepackage{fontspec}
\setmainfont{Arial}
\usepackage[margin=1.2cm,top=1.2cm,bottom=1.2cm]{geometry}
\usepackage{listings}

\setlength{\parindent}{0pt}
\setlength{\parskip}{4pt}

\lstset{
    basicstyle=\ttfamily\small,
    frame=single,
    breaklines=true,
    showstringspaces=false,
    tabsize=4
}

\begin{document}

% --- BAŞLIQ ---
\begin{center}
    {\large \textbf{BAKI ALİ NEFT MƏKTƏBİNİN NƏZDİNDƏ RƏQƏMSAL BİLİKLƏR LİSEYİ}}\\[0.1cm]
    {\Large \textbf{DƏRS 01 ÜZRƏ 10 DƏQİQƏLİK TƏKRAR QUİZ (RECAP QUIZ)}}\\[0.1cm]
    {\small \textbf{Fənn:} İnformatika \quad | \quad \textbf{Sinif:} 8-ci sinif \quad | \quad \textbf{Müddət:} 10 dəqiqə \quad | \quad \textbf{Maksimal Bal:} 20 bal}
\end{center}

\vspace{0.05cm}
\noindent\fbox{%
\parbox{0.98\textwidth}{%
\vspace{0.05cm}
\textbf{Şagirdin Adı, Soyadı:} \hrulefill \quad 
\textbf{Sinif:} \underline{\hspace{1.5cm}} \quad 
\textbf{Tarix:} \underline{\hspace{2cm}} \quad
\textbf{Toplanan Bal:} \underline{\hspace{1.2cm}} / 20
\vspace{0.05cm}
}}

\vspace{0.25cm}

\textbf{1. Test Sualı (4 bal)}\\
Şagirdin orta qiymətini (məsələn, \texttt{93.75}) yaddaşda saxlamaq üçün ən uyğun C++ məlumat tipi hansıdır?\\
$[ \ ]$ A) \texttt{int} \hspace{2cm} $[ \ ]$ B) \texttt{double} \hspace{2cm} $[ \ ]$ C) \texttt{string} \hspace{2cm} $[ \ ]$ D) \texttt{char}

\vspace{0.35cm}
\textbf{2. Test Sualı (4 bal)}\\
Aşağıdakı C++ kod sətrinin konsol çıxışı nə olacaq?
\begin{lstlisting}[language=C++]
cout << 17 % 5;
\end{lstlisting}
$[ \ ]$ A) \texttt{3.4} \hspace{2cm} $[ \ ]$ B) \texttt{3} \hspace{2cm} $[ \ ]$ C) \texttt{2} \hspace{2cm} $[ \ ]$ D) \texttt{0}

\vspace{0.35cm}
\textbf{3. Səhvi Tapın və Düzəldin (4 bal)}\\
Aşağıdakı C++ kod sətrində hansı kobud sintaksis xətası var? Düzəlişi yazın:
\begin{lstlisting}[language=C++]
cin << studentAge;
\end{lstlisting}
\textbf{Düzgün forma:} \hrulefill

\vspace{0.35cm}
\textbf{4. Çıxışı Əvvəlcədən Təyin Edin (Predict Output) (4 bal)}\\
Aşağıdakı kod parçası icra olunduqda ekrana hansı ədəd çıxacaq?

\begin{minipage}[t]{0.52\textwidth}
\begin{lstlisting}[language=C++]
int x = 10;
x += 6;
x--;
cout << x;
\end{lstlisting}
\end{minipage}
\hfill
\begin{minipage}[t]{0.44\textwidth}
\noindent\fbox{%
\parbox{0.95\linewidth}{%
\textbf{Konsol Çıxışı:}\\[0.2cm]
\vspace{0.7cm}
}}
\end{minipage}

\vspace{0.35cm}
\textbf{5. Qısa Kod Yazılışı (4 bal)}\\
Aşağıdakı 3 tələbi yerinə yetirən 3 sətirlik C++ kodunu yazın:
\begin{enumerate}
    \itemsep0em
    \item \texttt{radius} adında kəsr ədəd (\texttt{double}) dəyişəni elan edin.
    \item İstifadəçidən \texttt{cin} vasitəsilə bu radiusu daxil etməsini oxuyun.
    \item Çevrənin uzunluğunu (\texttt{2 * 3.14 * radius}) \texttt{cout} ilə ekrana çıxarın.
\end{enumerate}

\noindent\fbox{%
\parbox{0.98\textwidth}{%
\vspace{0.1cm}
\textbf{3 Sətirlik C++ Kodu:}\\[2.6cm]
}}

\vfill
\begin{center}
\footnotesize \textit{Müəllim: Avaz Əsgərov | BANM Nəzdində Rəqəmsal Biliklər Liseyi -- Səhifə 1 / 1}
\end{center}

\end{document}
